\section{Related Work}

\noindent\textbf{Transformer for Decision-Making.}~~
In general, reinforcement learning was proposed as a fundamentally online paradigm \cite{38}. The nature of online learning comes with some limitations when meeting the applications for which it is impossible to gather online data and learn simultaneously, such as autonomous driving. To this end, offline RL proposes that the agent can learn from a fixed dataset of previously collected data without gathering new data during learning \cite{39,40,41,42}. In the context of offline RL, recent works explored using transformer-based policy by treating RL tasks as a type of sequential prediction problem. Among them, a decision transformer \cite{2} is proposed to model trajectories as sequences and autoregressively predicts action conditioning on desired return-to-go, past states, and actions. Trajectory transformer \cite{43} demonstrated that transformer could learn single-task policies from offline data. Subsequently, the multi-game decision transformer \cite{14} and Gato \cite{27} further showed that transformer-based policies could address multi-tasks in the same domain and cross-domain tasks. However, these works focused on distilling expert policies from offline data and failed to enable self-improvement like IDT. Instead, when the offline data are sub-optimal, or the agent is required to adapt to new tasks, the multi-game decision transformers need to finetune the model parameters while Gato is required to get prompted with expert demonstrations.

\noindent\textbf{Meta RL.}~~
IDT falls into the category of methods of learning to learn, which is also known as meta-learning. More precisely, recent in-context RL methods can be categorized as in-context meta-RL methods. The general idea of learning self-improvement has a long history in RL but is limited to hyper-parameters in the early stages \cite{31}. In-context meta-RL methods \cite{33,32} are commonly trained in the online setting by maximizing multi-episodic value functions with memory-based architectures through environment interactions. Another online meta-RL attempts to find good network parameter initializations and then quickly adapt through additional gradient updates \cite{46,45}. More recently, meta-RL has seen substantial breakthroughs, from performance gains on popular benchmarks to offline settings, such as Bayesian RL \cite{36} and optimization-based meta-RL \cite{37}. Considering the difficulty of a completely offline setting, recent work has explored hybrid offline-online settings \cite{34,35}. IDT is similar to the hybrid offline-online setting, but the online phase does not involve gradient updates.

\noindent\textbf{In-Context RL.}~~
In-context RL is the one that addresses tasks by providing prompts or demonstrations \cite{2,43}. By training agents at a large scale, transformer-based policies usually have the ability to learn in context \cite{14,27}. The learning process is performed entirely in context and does not involve parameter updates of neural networks. In this work, we consider incremental in-context RL that involves learning from one's own behaviors through a trial-and-error manner. \citet{8} proposed Algorithm Distillation (AD) that automatically improves its performance in a trial-and-error manner by providing multiple historical trajectories. Subsequently, \citet{44} proposed a Decision-Pretrained Transformer, which trains the agent to find optimal behaviors faster by only predicting the optimal trajectory. More recently, \citet{1} further demonstrated that across-episodic contexts encourage large transformer models' emerging trial-and-error behaviors. However, these methods focus on the smallest action level, which causes across-episodic contexts to induce too-long sequences and suffer from huge computational costs. In contrast, IDT explores the trial-and-error ability of high-level decisions, which can significantly shorten the length of across-episodic contexts and, therefore, alleviate the computational costs arising from the self-attention mechanism.
